% !TeX program = xelatex
% UTF-8. Compile twice with XeLaTeX.
% Set \showsolutionsfalse below to create an exercise-only version.
\documentclass[11pt,a4paper,fontset=fandol]{ctexart}
\usepackage[margin=25mm]{geometry}
\usepackage{amsmath,amssymb,mathtools}
\usepackage{enumitem}
\usepackage[hidelinks]{hyperref}
\setlength{\parskip}{0.35em}
\setlist{itemsep=0.25em,topsep=0.4em}
\allowdisplaybreaks[2]
\newif\ifshowsolutions
\showsolutionstrue
\newcommand{\dd}{\mathrm{d}}
\newcommand{\R}{\mathbb{R}}
\title{引力研究微分几何短练习\\[0.4em]
\large 第01期：推前与拉回——从链式法则到诱导度规}
\author{}
\date{2026年10月28日}

\begin{document}
\maketitle
\noindent\textbf{学习目标：}
从映射与切向量的定义出发，计算一形式和度规的拉回，
并区分嵌入微分、时空度规、诱导度规的非退化性。
\textbf{前置知识：} 多元微积分、切向量与协变张量的基本含义、基础GR。
\textbf{练习时间：} 约30--45分钟；建议先做三道题，再看解答。

\section{核心概念}
设$\Phi:M\to N$是光滑映射，$p\in M$，$q=\Phi(p)$。
切向量$X_p\in T_pM$作用于函数，给出方向导数。
其\textbf{推前}是目标空间中的切向量，定义为
\begin{equation}
  (\Phi_{*p}X_p)[f]=X_p[f\circ\Phi],
  \qquad f\in C^\infty(N).
  \label{eq:push}
\end{equation}
映射将运动方向推向目标空间；目标空间中的测量规则可以拉回原空间。
一形式$\alpha_q\in T_q^*N$输入一个切向量并输出一个数。
其\textbf{拉回}定义为
\begin{equation}
  (\Phi^*\alpha)_p(X_p)
  =\alpha_{\Phi(p)}(\Phi_{*p}X_p).
  \label{eq:pull}
\end{equation}
因此，拉回的操作是：先把所有输入向量推前，再交给原来的张量。
对二阶协变张量（特别是度规）同样有
\begin{equation}
  (\Phi^*g)_p(X_p,Y_p)
  =g_{\Phi(p)}(\Phi_*X_p,\Phi_*Y_p).
  \label{eq:pullmetric}
\end{equation}
这些定义不要求$\Phi$可逆，也不要求流形已带有度规。
切向量总能逐点推前；但整个向量场未必能推成$N$上良好定义的向量场。
例如，多个源点映到同一目标点时，其向量的像可能不同。
一般情况下，逐点推前首先给出沿$\Phi$的向量场。
相关概念也为用流定义Lie导数作准备，参见文献\cite{tong2}。

\section{练习}
\subsection*{练习1：从定义到坐标公式}
源空间用坐标$u^a$，目标空间用坐标$x^\mu$，
映射写成$x^\mu=\Phi^\mu(u)$。设
\[
 X=X^a\partial_{u^a},\qquad
 \alpha=\alpha_\mu\,\dd x^\mu .
\]
重复指标求和，且函数微分定义为$\dd f(X)=X[f]$。
\begin{enumerate}[label=(\arabic*)]
  \item 从定义推导$\Phi_*X$和$\Phi^*\alpha$的坐标表达式。
  \item 核对配对恒等式
  $(\Phi^*\alpha)_p(X_p)=\alpha_{\Phi(p)}(\Phi_*X_p)$。
  \item 对任意$f\in C^\infty(N)$证明
  \[
    \boxed{\Phi^*(\dd f)=\dd(f\circ\Phi)}.
  \]
\end{enumerate}

\subsection*{练习2：一形式与面积形式的拉回}
取极坐标映射
\[
 \Phi(r,\varphi)=(r\cos\varphi,r\sin\varphi),
 \qquad r>0,
\]
其中$\varphi$限制在一个局部坐标区间内。
目标平面上的形式为
\[
 \alpha=-y\,\dd x+x\,\dd y,\qquad
 \omega=\dd x\wedge\dd y .
\]
\begin{enumerate}[label=(\arabic*)]
  \item 计算$\Phi_*\partial_r$与$\Phi_*\partial_\varphi$。
  \item 计算$\Phi^*\alpha$，并用配对关系检查结果。
  \item 计算$\Phi^*\omega$，解释径向因子的来源。
\end{enumerate}
所需楔积定义为
\[
 (\beta\wedge\gamma)(U,V)
 =\beta(U)\gamma(V)-\beta(V)\gamma(U).
\]
因此$\dd r\wedge\dd r=0$，
$\dd\varphi\wedge\dd r=-\dd r\wedge\dd\varphi$。

\subsection*{练习3：黑洞视界的诱导度规}
取$G=c=1$、质量参数$m>0$和度规号差$(-+++)$。
Schwarzschild时空在入射Eddington--Finkelstein坐标中的线元为
\begin{equation}
 \dd s^2=-F(r)\,\dd v^2+2\,\dd v\,\dd r
  +r^2(\dd\theta^2+\sin^2\theta\,\dd\varphi^2),
 \qquad F(r)=1-\frac{2m}{r}.
 \label{eq:ef}
\end{equation}
这组坐标在未来视界$r=2m$处保持正则，参见文献\cite{tong6}。
固定$R>0$，定义超曲面嵌入
\[
 \iota_R:\Sigma_R\hookrightarrow M,\qquad
 (v,\theta,\varphi)\mapsto(v,R,\theta,\varphi).
\]
计算可在$0<\theta<\pi$的球面坐标片内进行。
\begin{enumerate}[label=(\arabic*)]
  \item 计算$h=\iota_R^*g$。
  \item 判断$R>2m$与$R=2m$时$h$是否退化。
  \item 在$R=2m$处，嵌入的微分是否失去秩？时空度规是否退化？
  用计算解释诱导度规的结果为何不与它们矛盾。
\end{enumerate}

\ifshowsolutions
\clearpage
\section{详细解答}
\subsection{练习1}
对任意目标空间函数$f$，式\eqref{eq:push}给出
\begin{align}
 (\Phi_{*p}X_p)[f]
 &=X_p[f\circ\Phi]\notag\\
 &=X^a(p)\frac{\partial}{\partial u^a}f(\Phi(u))\bigg|_p\notag\\
 &=X^a(p)\frac{\partial\Phi^\mu}{\partial u^a}\bigg|_p
   \frac{\partial f}{\partial x^\mu}\bigg|_{\Phi(p)}.
\end{align}
最后一行是一个目标空间切向量作用于$f$，所以
\begin{equation}
 \boxed{\Phi_{*p}X_p
 =X^a(p)\,\partial_a\Phi^\mu(p)\,
   \partial_{x^\mu}\big|_{\Phi(p)}.}
\end{equation}
把它代入拉回定义：
\begin{align}
 (\Phi^*\alpha)_p(X_p)
 &=\alpha_{\Phi(p)}(\Phi_*X_p)\notag\\
 &=\alpha_\mu(\Phi(p))X^a(p)\partial_a\Phi^\mu(p)\notag\\
 &=\left[(\alpha_\mu\circ\Phi)\partial_a\Phi^\mu\right]_p X^a(p).
\end{align}
这对每个$X_p$都成立，而$\dd u^a(X_p)=X^a(p)$，故
\begin{equation}
 \boxed{\Phi^*\alpha
 =(\alpha_\mu\circ\Phi)\partial_a\Phi^\mu\,\dd u^a.}
\end{equation}
这一计算同时核对了配对恒等式：先拉回一形式再求值，
或先推前向量再求值，得到同一个数。

下面逐步证明函数微分的自然性：
\begin{align}
 [\Phi^*(\dd f)]_p(X_p)
 &=(\dd f)_{\Phi(p)}(\Phi_*X_p)
 &&\text{（拉回定义）}\notag\\
 &=(\Phi_*X_p)[f]
 &&\text{（函数微分定义）}\notag\\
 &=X_p[f\circ\Phi]
 &&\text{（推前定义）}\notag\\
 &=[\dd(f\circ\Phi)]_p(X_p).
\end{align}
两个一形式在任意点、任意切向量上的取值均相同，因而它们相等。
特别地，令$f=x^\mu$，便得到计算规则
\begin{equation}
 \boxed{\Phi^*(\dd x^\mu)=\dd\Phi^\mu.}
 \label{eq:basicrule}
\end{equation}

\subsection{练习2}
对映射的两个坐标函数求导，
\begin{align}
 \Phi_*\partial_r
 &=\cos\varphi\,\partial_x+\sin\varphi\,\partial_y,\\
 \Phi_*\partial_\varphi
 &=-r\sin\varphi\,\partial_x+r\cos\varphi\,\partial_y.
\end{align}
由式\eqref{eq:basicrule}，
\begin{align}
 \Phi^*(\dd x)
 &=\dd(r\cos\varphi)
 =\cos\varphi\,\dd r-r\sin\varphi\,\dd\varphi,\\
 \Phi^*(\dd y)
 &=\dd(r\sin\varphi)
 =\sin\varphi\,\dd r+r\cos\varphi\,\dd\varphi.
\end{align}
系数函数也要复合：
$\Phi^*x=r\cos\varphi$，$\Phi^*y=r\sin\varphi$。
因此
\begin{align}
 \Phi^*\alpha
 ={}&-r\sin\varphi
       (\cos\varphi\,\dd r-r\sin\varphi\,\dd\varphi)\notag\\
    &+r\cos\varphi
       (\sin\varphi\,\dd r+r\cos\varphi\,\dd\varphi)\notag\\
 ={}&(-r\sin\varphi\cos\varphi+r\cos\varphi\sin\varphi)\,\dd r\notag\\
    &+r^2(\sin^2\varphi+\cos^2\varphi)\,\dd\varphi\notag\\
 ={}&\boxed{r^2\,\dd\varphi}.
\end{align}
用$\partial_\varphi$作配对检查，左边为
$(\Phi^*\alpha)(\partial_\varphi)=r^2$；
在目标空间中则为
\begin{align}
 \alpha(\Phi_*\partial_\varphi)
 &=(-y\,\dd x+x\,\dd y)
  (-r\sin\varphi\,\partial_x+r\cos\varphi\,\partial_y)\notag\\
 &=r^2\sin^2\varphi+r^2\cos^2\varphi=r^2.
\end{align}

为计算面积形式，先从定义核对拉回与楔积的关系。
对任意$U,V\in T_pM$，
\begin{align}
 [\Phi^*(\beta\wedge\gamma)](U,V)
 &=(\beta\wedge\gamma)(\Phi_*U,\Phi_*V)\notag\\
 &=\beta(\Phi_*U)\gamma(\Phi_*V)
   -\beta(\Phi_*V)\gamma(\Phi_*U)\notag\\
 &=[(\Phi^*\beta)\wedge(\Phi^*\gamma)](U,V).
\end{align}
因此
\begin{align}
 \Phi^*(\dd x\wedge\dd y)
 &=(\Phi^*\dd x)\wedge(\Phi^*\dd y)\notag\\
 &=(\cos\varphi\,\dd r-r\sin\varphi\,\dd\varphi)
   \wedge(\sin\varphi\,\dd r+r\cos\varphi\,\dd\varphi)\notag\\
 &=r\cos^2\varphi\,\dd r\wedge\dd\varphi
   -r\sin^2\varphi\,\dd\varphi\wedge\dd r\notag\\
 &=\boxed{r\,\dd r\wedge\dd\varphi}.
\end{align}
其中径向因子来自映射的面积伸缩：
\begin{equation}
 \det\begin{pmatrix}
 \partial_r x&\partial_\varphi x\\
 \partial_r y&\partial_\varphi y
 \end{pmatrix}
 =
 \det\begin{pmatrix}
 \cos\varphi&-r\sin\varphi\\
 \sin\varphi&r\cos\varphi
 \end{pmatrix}=r.
\end{equation}
这给出了极坐标积分中径向因子的几何来源。

\subsection{练习3}
先逐项拉回坐标函数：
\[
 v\circ\iota_R=v,\qquad r\circ\iota_R=R.
\]
$R$是$\Sigma_R$上的常数，故
\[
 \iota_R^*(\dd v)=\dd v,\qquad
 \boxed{\iota_R^*(\dd r)=\dd R=0}.
\]
另外，
$\iota_R^*(\dd\theta)=\dd\theta$，
$\iota_R^*(\dd\varphi)=\dd\varphi$。
所以
\begin{equation}
 \boxed{h=-F(R)\,\dd v^2
       +R^2\,\dd\theta^2+R^2\sin^2\theta\,\dd\varphi^2.}
 \label{eq:induced}
\end{equation}
这里$\dd v^2$表示$\dd v\otimes\dd v$；
线元中的$2\,\dd v\,\dd r$表示
$\dd v\otimes\dd r+\dd r\otimes\dd v$。
度规使用对称张量积，而不是楔积。

当$R>2m$时，$F(R)>0$，且
\[
 (h_{ab})=\operatorname{diag}(-F(R),R^2,R^2\sin^2\theta).
\]
它有一个负方向和两个正方向；在所选坐标片中，
\[
 \det h=-F(R)R^4\sin^2\theta\ne0.
\]
因此$\Sigma_R$是类时超曲面，其诱导度规非退化。

当$R=2m$时，
\[
 h=(2m)^2(\dd\theta^2+\sin^2\theta\,\dd\varphi^2).
\]
对任意
$Y=A\partial_v+B\partial_\theta+C\partial_\varphi\in T\Sigma_{2m}$，
都有
\[
 h(\partial_v,Y)=0.
\]
非零向量$\partial_v$属于$h$的核，故诱导度规退化。

另一方面，嵌入的微分仍满足
\[
 \iota_{R*}\partial_v=\partial_v,\qquad
 \iota_{R*}\partial_\theta=\partial_\theta,\qquad
 \iota_{R*}\partial_\varphi=\partial_\varphi.
\]
其像仍有三个线性无关方向，秩为$3$。
环境时空度规的$v,r$子块为
\[
 \begin{pmatrix}-F(r)&1\\1&0\end{pmatrix},
 \qquad \det=-1,
\]
所以在视界处依然可逆。
角向子块也在所选坐标片内非退化，故完整时空度规非退化。
特别地，
\[
 g(\partial_v,\partial_r)=1.
\]
$\partial_v$只与视界上的所有切向方向正交，
没有与全部时空方向正交；$\partial_r$不属于$T\Sigma_{2m}$，
因为$\dd r(\partial_r)=1$而切向量必须满足$\dd r(Y)=0$。

还可直接核对零超曲面的法向量。
$v,r$子块的逆矩阵为
\[
 \begin{pmatrix}0&1\\1&F(r)\end{pmatrix}.
\]
因而法向余向量$n=\dd r$的度规对偶为
\[
 n^\sharp=g^{\mu r}\partial_\mu=\partial_v+F(r)\partial_r.
\]
在视界上$F=0$，于是$n^\sharp=\partial_v$同时是切向量，
且对任何切向量$Y$，
$g(n^\sharp,Y)=n(Y)=\dd r(Y)=0$。
这解释了诱导度规为何存在非零核。

需要区别两种性质：
\[
 \underbrace{h(X,Y)=0\ \text{对所有 }Y}_{\text{非零核向量，表明退化}}
 \qquad\text{与}\qquad
 \underbrace{h(X,X)=0}_{\text{零向量，不足以表明退化}},
 \quad X\ne0.
\]
非退化的Lorentz度规可以含有零向量。
本题视界诱导度规满足前一种更强的性质。

\section{引力应用与工具卡}
\textbf{具体应用：}
计算有限半径边界、黑洞视界或嵌入曲面的内禀几何时，
首先拉回时空度规。
对于零超曲面，不能假定诱导三维度规存在普通逆矩阵；
视界的二维球面截面仍有非退化度规，
沿生成线的方向需要单独处理。
这种区分是后续研究零边界、视界生成元和边界荷的基础。

\begin{enumerate}[label=\textbf{\arabic*.}]
 \item \textbf{拉回一形式：}
 系数函数先复合，基底使用$\Phi^*(\dd x^\mu)=\dd\Phi^\mu$。
 \item \textbf{拉回度规：}
 若嵌入写成$x^\mu=X^\mu(y^a)$，则
 \[
 h_{ab}=g_{\mu\nu}(X(y))
 \frac{\partial X^\mu}{\partial y^a}
 \frac{\partial X^\nu}{\partial y^b}.
 \]
 该式来自把$\partial_{y^a}$与$\partial_{y^b}$分别推前后代入$g$。
 \item \textbf{分别检查三个对象：}
 映射微分的秩、环境度规的非退化性、诱导度规的非退化性。
 它们是不同的问题。
\end{enumerate}

\textbf{与后续单元的联系：}
本期的$\Phi^*(\dd f)=\dd(\Phi^*f)$将在外微分单元中推广到任意微分形式。
之后令$\Phi$成为向量场生成的一参数流，便可以通过拉回的变化率定义Lie导数。
本期不把这一预告当作已证明的结论。

\textbf{下一单元：}
楔积与外微分的实际运算，特别是$r\,\dd\varphi$这样的非坐标余标架，
为Cartan结构方程准备计算基础。
\fi

\begin{thebibliography}{9}
 \bibitem{tong2}
 David Tong, \emph{General Relativity},
 Chapter 2: Introducing Differential Geometry.
 \url{https://www.davidtong.org/teaching/general-relativity/grhtml/S2}.
 \bibitem{tong6}
 David Tong, \emph{General Relativity},
 Section 6.1.3: Eddington--Finkelstein Coordinates.
 \url{https://davidtong.org/teaching/general-relativity/grhtml/S6}.
\end{thebibliography}

\noindent\textbf{系列记录：}
第01期；核心技能为推前、拉回、诱导度规；
题目指纹为链式法则、极坐标形式拉回、Schwarzschild视界诱导度规。
本期材料已提供，用户完成与掌握状态待反馈。
下一期为第02单元：楔积与外微分的实际运算。

\end{document}

